\section{Conceptos fundamentales}

\textbf{Exactitud:} Cercanía entre el valor medido y el valor verdadero o de referencia.

\textbf{Precisión:} Grado de concordancia entre mediciones repetidas de una misma muestra (repetitividad o reproducibilidad).

\textbf{Incertidumbre:} Parámetro asociado al resultado de una medición que caracteriza la dispersión de los valores que podrían ser atribuidos al mensurando.

\textbf{Material de referencia certificado (MRC):} Material que contiene una concentración certificada de un analito, utilizado para verificar la exactitud.

\textbf{Trazabilidad:} Capacidad de relacionar un resultado de medición con una referencia a través de una cadena ininterrumpida de comparaciones.


\section{Tipos de errores}
\begin{itemize}
    \item \textbf{Errores sistemáticos (determinados):} Afectan la exactitud. Pueden ser del método, del instrumento o del analista. Son reproducibles y se pueden corregir.
    \item \textbf{Errores aleatorios (indeterminados):} Afectan la precisión. Son debidos a variaciones inevitables en el proceso de medición.
\end{itemize}

\section{Curvas de calibración y validación}
\begin{enumerate}
    \item \textbf{Construcción de la curva:} Se preparan estándares de concentración conocida y se mide su señal instrumental.
    \item \textbf{Linealidad:} La curva debe ser lineal en el rango de trabajo.
    \item \textbf{Límite de detección (LOD) y cuantificación (LOQ):} Se determinan a partir de la curva de calibración.
    \item \textbf{Estudios de repetitividad y reproducibilidad (R\&R):} Evalúan la precisión del método.
\end{enumerate}

\resumebox{dato}{Fuerza magnética del protón}{\footnotesize
    \textbf{Normas aplicables:} En México, las normas NOM (Normas Oficiales Mexicanas) y las normas ASTM (American Society for Testing and Materials) son de referencia para los métodos analíticos.
}

\resumebox{ejercicio}{Cálculo de parámetros de validación}{\footnotesize
    Se obtuvieron los siguientes datos de una curva de calibración para la determinación de glucosa: $y = 0.05x + 0.001$ ($R^2=0.999$). ¿Cuál es la concentración de una muestra con una absorbancia de 0.5? ¿Cómo estimarías el LOD y LOQ?
}